\documentclass[10pt]{beamer}
\usetheme{metropolis}
\usepackage{graphicx}
\usepackage{listings}
\usepackage{booktabs}
\usepackage{xcolor}
\definecolor{codebg}{HTML}{F5F5F5}
\definecolor{codegreen}{HTML}{2E7D32}
\definecolor{codegray}{HTML}{757575}
\definecolor{codeblue}{HTML}{1565C0}
\lstdefinestyle{rstyle}{language=R,backgroundcolor=\color{codebg},basicstyle=\ttfamily\scriptsize,
  keywordstyle=\color{codeblue}\bfseries,stringstyle=\color{codegreen},
  commentstyle=\color{codegray}\itshape,breaklines=true,frame=single,
  rulecolor=\color{codegray!30},showstringspaces=false,tabsize=2,
  morekeywords={library,ggplot,aes,geom_histogram,geom_boxplot,geom_point,geom_qq,
    geom_qq_line,summary,glimpse,str,head,tail,dim,nrow,ncol,
    theme_minimal,labs,skimr,skim,DataExplorer,create_report}}
\lstdefinestyle{pythonstyle}{language=Python,backgroundcolor=\color{codebg},basicstyle=\ttfamily\scriptsize,
  keywordstyle=\color{codeblue}\bfseries,stringstyle=\color{codegreen},
  commentstyle=\color{codegray}\itshape,breaklines=true,frame=single,
  rulecolor=\color{codegray!30},showstringspaces=false,tabsize=4,
  morekeywords={import,from,as,pd,plt,sns,np,fig,ax,describe,info,shape,head,
    probplot,hist,boxplot,violinplot}}
\title{Tukey's EDA Philosophy}
\subtitle{Workshop 4 --- Module 1}
\author{Prof.Asoc.\ Endri Raco}
\institute{Data Visualization for Data Scientists \\ UNYT Tirana}
\date{2025/26}
\begin{document}

\begin{frame}
  \titlepage
\end{frame}

\begin{frame}{Module Roadmap}
  \begin{enumerate}
    \item Confirmatory vs exploratory: two cultures of data analysis
    \item Tukey's detective metaphor (1977)
    \item The original EDA toolkit: stem-and-leaf, boxplot, residuals
    \item Letter-value plots: extending the box for large $n$
    \item The 4-plot diagnostic panel
    \item The modern EDA workflow
  \end{enumerate}
  \begin{exampleblock}{Learning Outcome}
    You will understand EDA as a philosophy (not a technique), adopt
    Tukey's iterative ``look at the data first'' mindset, and apply
    the 4-phase EDA workflow to any new dataset.
  \end{exampleblock}
\end{frame}

\section{Two Cultures}

\begin{frame}{Confirmatory vs Exploratory}
  \begin{center}
    \includegraphics[width=0.88\textwidth]{figures_w04m01/cda_vs_eda}
  \end{center}
  \begin{alertblock}{Pro-Tip}
    Most statistics courses teach CDA: hypothesis $\rightarrow$ test $\rightarrow$
    $p$-value. Tukey argued that this skips the most important step:
    \textbf{looking at the data}. EDA generates the hypotheses that CDA
    then tests. Both are necessary; EDA comes first.
  \end{alertblock}
\end{frame}

\begin{frame}{Tukey's Own Words}
  \begin{quote}
    ``Exploratory data analysis is detective work --- numerical detective
    work --- or counting detective work --- or graphical detective work.''
  \end{quote}
  \hfill --- John W. Tukey, \emph{Exploratory Data Analysis} (1977)
  \vspace{8pt}

  \begin{quote}
    ``The greatest value of a picture is when it forces us to notice
    what we never expected to see.''
  \end{quote}
  \hfill --- John W. Tukey
  \vspace{8pt}

  \begin{exampleblock}{Data Insight}
    Tukey invented the boxplot, the stem-and-leaf display, the term
    ``bit'' (with Shannon), and the FFT algorithm (with Cooley). His
    1977 book \emph{Exploratory Data Analysis} founded the field.
  \end{exampleblock}
\end{frame}

\section{The Detective Metaphor}

\begin{frame}{EDA as Investigation: Five Steps}
  \begin{center}
    \includegraphics[width=0.92\textwidth]{figures_w04m01/detective}
  \end{center}
\end{frame}

\begin{frame}{Key Principles of EDA}
  \centering
  \begin{tabular}{lp{6cm}}
    \toprule
    \textbf{Principle} & \textbf{Implication} \\
    \midrule
    Scepticism & Assume nothing about the data; let it speak \\
    Iteration & One plot leads to the next question \\
    Robustness & Use medians over means; resist outlier influence \\
    Residuals & After fitting a pattern, examine what remains \\
    Graphical primacy & Tables summarise; plots reveal \\
    Re-expression & Transform to simplify (log, sqrt, reciprocal) \\
    \bottomrule
  \end{tabular}
\end{frame}

\section{Tukey's Original Toolkit}

\begin{frame}{Stem-and-Leaf: Every Value + Shape}
  \begin{center}
    \includegraphics[width=0.55\textwidth]{figures_w04m01/stem_leaf}
  \end{center}
  \begin{exampleblock}{Data Insight}
    The stem-and-leaf preserves every data value (unlike a histogram which
    bins) while showing the distribution shape. Modern equivalent:
    \texttt{stem()} in R, or simply a high-resolution histogram. Useful
    for $n < 100$ in teaching contexts.
  \end{exampleblock}
\end{frame}

\begin{frame}{From Tukey (1977) to Modern R/Python}
  \begin{center}
    \includegraphics[width=0.92\textwidth]{figures_w04m01/tukey_toolkit}
  \end{center}
\end{frame}

\begin{frame}{Letter-Value Plot: Beyond the 5-Number Summary}
  \begin{center}
    \includegraphics[width=0.88\textwidth]{figures_w04m01/letter_value}
  \end{center}
  \begin{alertblock}{Pro-Tip}
    For large $n$ ($> 200$), the standard boxplot wastes information:
    it shows only 5 quantiles. The \textbf{letter-value plot} (Hofmann
    et al., 2017) adds nested boxes at the eighths (E), sixteenths (D),
    etc., revealing tail behaviour. In R: \texttt{ggplot2::geom\_lv()}.
    In Python: \texttt{sns.boxenplot()}.
  \end{alertblock}
\end{frame}

\section{The 4-Plot}

\begin{frame}{Tukey's 4-Plot Diagnostic Panel}
  \begin{center}
    \includegraphics[width=0.82\textwidth]{figures_w04m01/four_plot}
  \end{center}
\end{frame}

\begin{frame}{What Each Panel Reveals}
  \centering
  \begin{tabular}{lp{5.5cm}}
    \toprule
    \textbf{Panel} & \textbf{What It Checks} \\
    \midrule
    (a) Run sequence & Are there shifts, trends, or cycles over time/index? \\
    (b) Lag-1 plot & Is there autocorrelation? (If cloud $\approx$ circle: no) \\
    (c) Histogram & What is the shape? Unimodal, skewed, bimodal? \\
    (d) Normal Q-Q & Do the data follow a normal distribution? Departures at tails? \\
    \bottomrule
  \end{tabular}
  \vspace{6pt}

  {\small If (a) shows no trend, (b) shows no autocorrelation, (c) looks
  unimodal, and (d) follows the diagonal --- the data are well-behaved
  for parametric analysis. Deviations guide further EDA.}
\end{frame}

\begin{frame}[fragile]{The 4-Plot in R vs Python}
  \begin{columns}[T]
    \begin{column}{0.48\textwidth}
      \textbf{R}
\begin{lstlisting}[style=rstyle]
library(patchwork)

p_run <- ggplot(df, aes(x = 1:n(),
    y = value)) + geom_line() +
  geom_hline(yintercept = mean(
    df$value), color = "#E53935",
    linetype = "dashed") +
  labs(title = "(a) Run Sequence")

p_lag <- ggplot(df, aes(
    x = lag(value), y = value)) +
  geom_point(alpha = 0.3) +
  labs(title = "(b) Lag-1 Plot")

p_hist <- ggplot(df, aes(
    x = value)) +
  geom_histogram(bins = 25) +
  labs(title = "(c) Histogram")

p_qq <- ggplot(df, aes(
    sample = value)) +
  geom_qq() + geom_qq_line(
    color = "#E53935") +
  labs(title = "(d) Q-Q Plot")

(p_run | p_lag) / (p_hist | p_qq)
\end{lstlisting}
    \end{column}
    \begin{column}{0.48\textwidth}
      \textbf{Python}
\begin{lstlisting}[style=pythonstyle]
from scipy.stats import probplot

fig, axes = plt.subplots(2, 2)

# (a) Run sequence
axes[0,0].plot(range(n), values)
axes[0,0].axhline(
    y=np.mean(values),
    color="#E53935", ls="--")
axes[0,0].set_title("(a) Run")

# (b) Lag-1
axes[0,1].scatter(
    values[:-1], values[1:],
    s=8, alpha=0.3)
axes[0,1].set_title("(b) Lag-1")

# (c) Histogram
axes[1,0].hist(values, bins=25)
axes[1,0].set_title("(c) Hist")

# (d) Q-Q
probplot(values, dist="norm",
    plot=axes[1,1])
axes[1,1].set_title("(d) Q-Q")
\end{lstlisting}
    \end{column}
  \end{columns}
\end{frame}

\section{Modern EDA Workflow}

\begin{frame}{The Four-Phase EDA Workflow}
  \begin{center}
    \includegraphics[width=0.92\textwidth]{figures_w04m01/eda_workflow}
  \end{center}
\end{frame}

\begin{frame}[fragile]{Phase 1: Import \& Inspect}
  \begin{columns}[T]
    \begin{column}{0.48\textwidth}
      \textbf{R}
\begin{lstlisting}[style=rstyle]
library(tidyverse)

df <- read_csv("data.csv")
dim(df)           # rows x cols
glimpse(df)       # types + preview
summary(df)       # 5-num + NA count

# Comprehensive EDA report
library(skimr)
skim(df)           # grouped stats

# Or auto-generated report
library(DataExplorer)
create_report(df)  # HTML report
\end{lstlisting}
    \end{column}
    \begin{column}{0.48\textwidth}
      \textbf{Python}
\begin{lstlisting}[style=pythonstyle]
import pandas as pd

df = pd.read_csv("data.csv")
df.shape           # (rows, cols)
df.info()          # types + non-null
df.describe()      # stats for numeric
df.head()          # first 5 rows

# Comprehensive EDA
# pip install ydata-profiling
from ydata_profiling import (
    ProfileReport)
report = ProfileReport(df)
report.to_file("eda_report.html")
\end{lstlisting}
    \end{column}
  \end{columns}
\end{frame}

\section{Complete Examples}

\begin{frame}[fragile]{R: EDA First Look}
  \begin{columns}[T]
    \begin{column}{0.52\textwidth}
\begin{lstlisting}[style=rstyle]
library(tidyverse); library(patchwork)
apps <- read_csv(
  "googleplaystore.csv") |>
  filter(!is.na(Rating))

p1 <- ggplot(apps, aes(x=Rating))+
  geom_histogram(bins=30,
    fill="#1565C0", color="white") +
  theme_minimal() +
  labs(title="(a) Histogram")

p2 <- ggplot(apps, aes(x="",
    y=Rating)) +
  geom_boxplot(fill="#BBDEFB",
    width=0.3) +
  theme_minimal() +
  labs(title="(b) Boxplot", x=NULL)

p3 <- ggplot(apps,
    aes(sample=Rating)) +
  geom_qq(alpha=0.1, size=0.5) +
  geom_qq_line(color="#E53935") +
  theme_minimal() +
  labs(title="(c) Q-Q Plot")

p1 | p2 | p3
\end{lstlisting}
    \end{column}
    \begin{column}{0.45\textwidth}
      \includegraphics[width=\textwidth]{figures_w04m01/r_output}
    \end{column}
  \end{columns}
\end{frame}

\begin{frame}[fragile]{Python: EDA First Look}
  \begin{columns}[T]
    \begin{column}{0.52\textwidth}
\begin{lstlisting}[style=pythonstyle]
import pandas as pd
import matplotlib.pyplot as plt
from scipy.stats import probplot

apps = pd.read_csv(
    "googleplaystore.csv")
ratings = apps["Rating"].dropna()

fig, axes = plt.subplots(1, 3,
    figsize=(10, 3.5))

axes[0].hist(ratings, bins=30,
    color="#2E7D32",
    edgecolor="white", lw=0.3)
axes[0].set_title("Histogram")

parts = axes[1].violinplot(
    [ratings],
    showmeans=True,
    showmedians=True)
axes[1].set_title("Violin")

probplot(ratings, dist="norm",
    plot=axes[2])
axes[2].set_title("Q-Q Plot")
\end{lstlisting}
    \end{column}
    \begin{column}{0.45\textwidth}
      \includegraphics[width=\textwidth]{figures_w04m01/py_output}
    \end{column}
  \end{columns}
\end{frame}

\section{Synthesis}

\begin{frame}{Key Takeaways}
  \begin{enumerate}
    \item \textbf{EDA is a philosophy}, not a technique: look at the data
          before modelling. CDA tests hypotheses; EDA generates them.
    \item Tukey's \textbf{detective metaphor}: investigate, form hypotheses,
          check residuals, iterate. Never trust a single plot.
    \item The \textbf{4-plot} (run sequence + lag-1 + histogram + Q-Q) is
          the first diagnostic panel for any univariate analysis.
    \item \textbf{Letter-value plots} extend boxplots for large $n$ by
          showing nested quantile boxes: \texttt{geom\_lv()} or
          \texttt{sns.boxenplot()}.
    \item The \textbf{modern EDA workflow}: Import/Inspect $\rightarrow$
          Clean/Transform $\rightarrow$ Visualize/Explore $\rightarrow$
          Document/Communicate.
    \item Automate Phase 1 with \texttt{skimr::skim()} (R) or
          \texttt{ydata-profiling} (Python) for instant EDA reports.
  \end{enumerate}
\end{frame}

\begin{frame}{Essential Readings}
  \begin{itemize}
    \item Tukey, J.~W. (1977). \emph{Exploratory Data Analysis}.
          Addison-Wesley. (The founding text.)
    \item Hofmann, H., Wickham, H. \& Kafadar, K. (2017). ``Letter-value
          plots.'' \emph{Journal of Computational and Graphical Statistics}.
    \item Wickham, H. \& Grolemund, G. (2023). \emph{R for Data Science},
          2nd ed., Chapters 2, 10 (Workflow, EDA).
    \item Wilke, C.~O. (2019). \emph{Fundamentals of Data Visualization},
          Chapter 7 (Histograms and Density).
  \end{itemize}
  \vspace{6pt}
  \textbf{Next Module}: Data Wrangling for Visualization in R (W04-M02)
\end{frame}

\end{document}
